% Nouveau EGA V (cote 25): archive cote 25, site page 160 (PDF page 161).
% typescript. Draft compared: none.
% AI-transcribed by Claude Opus 5.5 (claude-opus-5-5, adaptive thinking, effort high), Message Batches API, batch ega25b (msgbatch_01A7n5YuXubUPNK1zsepi8RG), 2026-11-13; not yet checked by eye.
\scanpage{25/160}
\mspage{\hand{3}}

proposition semble correct , et devrait se prouver de façon toute analogue
en montrant que pour $x\in X$\inserted{$^{(k)}$} donné, il y a une hypersurface \inserted{$H$} de degré $n\geqslant 2$
donné qui est tangente à X en x, et telle que x soit un point quadratique
élémentaire de $X$\hand{\unsure{$\cap$}}$H$ , de dimension une de moins .... Bien mettre les
pieds dans le plat, en faisant remarquer que si k \uline{est} de car.2 et X
connexe \inserted{non vide} de dimension impaire, alors les conditions de 8.8. \uline{ne sont pas}
vérifiées i.e. L/K est nécessairement inséparable (je crois, de degré
2 exactement si $X\subset P^r$ , mais sans garantie ...).

15.1. Soient S un préschéma, $\underline{E}$ un Module loc libre de type fini
sur S, $\phi$ une section de $\underline{\mathrm{Sym\xout{m}}}^n(\underline{E})$ , i.e. une ``n-forme sur $\check{\underline{E}}$ '' , où $n$\over{$>$}{\unsure{$\geqslant$}}$0$.
La donnée de $\phi$ équivaut à la donnée d'une section de $\underline{O}_P(n)$ sur $P=P(\underline{E})$
(réf àIII), et définit donc un sous-préschéma $V(\phi)$ de P. Nous dirons que
la forme $\phi$ est \uline{lisse} si $V(\phi)$ est lisse, et si de plus pour tout $s\in S$,
$V(\phi)_s\neq P_s$ (i.e. $\phi(s)\in \mathrm{Sym}^n_{k(s)}E\otimes k(s)$ est \over{$=$}{\unsure{$\neq$}}$0$ ). On voit aussitôt
que $\phi$ est lisse sss pour tout \xout{x} $\in$S, $\phi(s)$ est lisse, \xout{\illegible{3}}
\xout{\unsure{d'un corps}}(ce qui nous ramène au cas où S est le spectre d'un corps) et
que la notion de forme lisse est invariante par changement de corps de
base (ce quin nous ramène en particulier au cas d'un corps de base alg
clos); on peut résumer ces deux propriétés en disant que si $S'\to S$ est
un morphisme surjectif, alors $\phi$ est lisse sss son image inverse $\phi'$ sur S'
l'est. Bien entendu, si E est un module projectif de type fini sur un
anneau A, et $\phi\in\mathrm{Sym}^n_A(E)$, on dit que $\phi$ est lisse si .... Lorsque $E=A^{\unsure{N}+1}$,
la donnée de $\phi$ équivaut à celle d'un\xout{e forme} polynôme homogène de degré n
en des variables $X_0,..,X_r$ , et le critère jacobien implique immédiatement
que \inserted{\unsure{$n$ invers.}} $\phi$ est lisse sss l'idéal engendré par $\phi$ et les $\dfrac{\partial\phi}{\partial X_i}$ contient
une puissance de l'idéal d'augmentation $(X_0,\dots,X_r)$ , i.e. contient une

\margin{160} \margin{\unsure{Car mais on} trouve les sections inversibles de $\underline{O}_S$ !}
